\documentclass{article}

\usepackage[utf8]{inputenc}     % Soporte para tildes y eñes
\usepackage[T1]{fontenc}
\usepackage[spanish]{babel}     % Para usar títulos en español si querés
\usepackage{amsmath, amssymb}   % Símbolos matemáticos
\usepackage{xcolor}             % Colores para código
\usepackage{listings}           % Código fuente
\usepackage{geometry}           % Márgenes
\usepackage{array}


\geometry{margin=2.5cm}


% Configuración para mostrar C++ lindo
\lstdefinestyle{codestyle}{
	basicstyle=\scriptsize\ttfamily,
    keywordstyle=\bfseries,
    emph={int,char,double,float,unsigned,void,bool},
    emphstyle=\bfseries,
    breakatwhitespace=false,
    breaklines=true,
    captionpos=b,
    keepspaces=false,
    numbers=left,
    numberstyle=\tiny\ttfamily,
    numbersep=5pt,
    frame=single,
    framesep=3pt,
    showspaces=false,
    showstringspaces=false,
    showtabs=false,
    tabsize=2,
    commentstyle=\color{commentgreen}, % comment color
    keywordstyle=\color{blue}, % keyword color
    stringstyle=\color{red} % string color
}
\lstset{style=codestyle}

\title{Random}
\author{ChichuMerge}
\date{July 30, 2025}

\begin{document}

\maketitle

\section{ST Lechugero?}
\begin{lstlisting}
//El ST era de minimo
//Vos arrancabas en un arbol y corrias hacia su izquierda queres saber cuando ves 2 veces mismo arbol o llegas a 0
unordered_map<int, set<int>> v;    
vector<int> arboles(n);
forn(i, n){
  int x; cin >> x;
  arboles[i] = x;
  v[x].insert(i); //para cada tipo llevas un set con los indices
}

ST st(n);

for(auto &it : v){ //recorro el mapa
  for(auto pos = it.sc.begin(); pos != it.sc.end(); pos++){ //recorro el set
    int i = *pos;
    auto sig = next(pos);
    st[i] = sig != it.sc.end() ? *sig : INF; //actualizo cual es el sig del arbol i
  }
}
st.updall();

while(q--){
    char c; cin >> c;

    if(c == 'E'){
        int x; cin >> x; x--;
        int l = -1, r = x;

        //binaria para hallar la respuesta
        while(r - l > 1){
        int m = (r + l) / 2;
        int val = st.get(m, x); //encuentro el sig mas chico
        if(val <= x) l = m; 
        else r = m;
        }
        cout << (l == -1 ? x + 1 : x - l) << '\n'; //llego al rio, distancia de yo al arbol

    }else{
        //actualizacion horrible de sig
        int x, y; cin >> x >> y;
        x--;
        int pos = arboles[x]; //tipo del arbol a sacar
        auto act = v[pos].lower_bound(x); //apunto al que voy a sacar
        if(act != v[pos].begin()){
        auto sig = next(act);
        auto ant = prev(act);
        if(sig == v[pos].end()){
            st.set(*ant, INF);
        }else st.set(*ant, *sig);
        }
        v[pos].erase(x);

        arboles[x] = y;
        v[y].insert(x); //inserto el nuevo arbol en el tipo correspondiente
        act = v[y].lower_bound(x); //lo busco

        if(act != v[y].begin()){
        auto ant = prev(act);
        st.set(*ant, x);
        }
        auto sig = next(act);
        st.set(x, (sig == v[y].end() ? INF : *sig));
    }
}
\end{lstlisting}


\section{Nunca sabes cuando necesitas un triangulo de Pascal}
\begin{center}
\begin{tabular}{>{$}l<{$}|*{15}{c}}
\multicolumn{1}{l}{$n$}                     \\\cline{1-1}
0  & 1    &      &      &      &      &      &      &      &      &      &      &      &      &      &      \\
1  & 1    & 1    &      &      &      &      &      &      &      &      &      &      &      &      &      \\
2  & 1    & 2    & 1    &      &      &      &      &      &      &      &      &      &      &      &      \\
3  & 1    & 3    & 3    & 1    &      &      &      &      &      &      &      &      &      &      &      \\
4  & 1    & 4    & 6    & 4    & 1    &      &      &      &      &      &      &      &      &      &      \\
5  & 1    & 5    & 10   & 10   & 5    & 1    &      &      &      &      &      &      &      &      &      \\
6  & 1    & 6    & 15   & 20   & 15   & 6    & 1    &      &      &      &      &      &      &      &      \\
7  & 1    & 7    & 21   & 35   & 35   & 21   & 7    & 1    &      &      &      &      &      &      &      \\
8  & 1    & 8    & 28   & 56   & 70   & 56   & 28   & 8    & 1    &      &      &      &      &      &      \\
9  & 1    & 9    & 36   & 84   & 126  & 126  & 84   & 36   & 9    & 1    &      &      &      &      &      \\
10 & 1    & 10   & 45   & 120  & 210  & 252  & 210  & 120  & 45   & 10   & 1    &      &      &      &      \\
11 & 1    & 11   & 55   & 165  & 330  & 462  & 462  & 330  & 165  & 55   & 11   & 1    &      &      &      \\
12 & 1    & 12   & 66   & 220  & 495  & 792  & 924  & 792  & 495  & 220  & 66   & 12   & 1    &      &      \\
13 & 1    & 13   & 78   & 286  & 715  & 1287 & 1716 & 1716 & 1287 & 715  & 286  & 78   & 13   & 1    &      \\
14 & 1    & 14   & 91   & 364  & 1001 & 2002 & 3003 & 3432 & 3003 & 2002 & 1001 & 364  & 91   & 14   & 1    \\
\hline
\multicolumn{1}{l}{} 
   & 0 & 1 & 2 & 3 & 4 & 5 & 6 & 7 & 8 & 9 & 10 & 11 & 12 & 13 & 14 \\\cline{2-16}
\multicolumn{1}{l}{} &\multicolumn{15}{c}{$k$}
\end{tabular}
\end{center}

\end{document}
